This thesis evaluates the impact of generative artificial intelligence on video game development by analyzing the differences in player experience and perceived quality between traditionally programmed code and AI-assisted code in the Godot Engine. The study addresses the question of whether the integration of generative models into software engineering degrades the functional quality of the product or whether variations in its evaluation stem from cognitive biases regarding the technology’s authorship.

To this end, an A/B experimental design was implemented with a multidisciplinary sample of 60 college students, who evaluated two-dimensional prototypes using the standardized Player Experience Inventory (PXI) questionnaire under both blind and informed conditions. The quantitative results demonstrated functional and technical parity in the blind test, with no significant differences in the gaming experience. However, explicitly revealing the authorship to the informed group led to lower ratings for the AI-generated version.

It is concluded that AI-generated code achieves performance equivalent to that of human development, but its perception is significantly influenced by labeling bias and the user’s prior attitudes, which allowed the research hypothesis to be accepted. It is recommended that the study be extended to specialized audiences and commercial industries, that longitudinal evaluations be implemented, and that testing conditions involving partial information be incorporated to distinguish actual detection capabilities from induced bias.